Let $(S_j)_{j\ge0}$ be simple random walk on $\mathbb Z^2$ starting at the origin. For each realized path let $R_N$ be the largest integer $r\ge0$ such that every lattice point in the Euclidean disk of radius $r$ has been visited by time $N$. This is a random radius. The literal alternative requiring coverage by time $N$ on almost every path gives no positive radius: the event of taking $N$ consecutive steps in one direction has positive probability and misses a neighbor of the origin.

With the intended pathwise definition the precise distributional conclusion is
\[
 \frac{(\log(R_N\vee1))^2}{\log N}\ \Longrightarrow\ \operatorname{Exp}(4),\tag{1}
\]
where the exponential distribution has \emph{survival} function $e^{-4x}$. Thus the editorial's displayed $e^{-4x}$ cannot be the cumulative distribution function written there: it is decreasing. We prove the inversion and its scale consequences carefully, using the deep cover-time theorem of Dembo--Peres--Rosen--Zeitouni, \emph{Cover times for Brownian motion and random walks in two dimensions}, Theorem 1.4, \url{https://annals.math.princeton.edu/wp-content/uploads/annals-v160-n2-p02.pdf}, as an explicit external input.
